\begin{abcliste}
\abc
Ein Quader kippt um seine untere (talseitige) Kante. Die Gewichtskraft greift im Schwerpunkt S an und zeigt senkrecht nach unten. Verläuft ihre Wirkungslinie innerhalb der Kante (unter der Grundfläche), dreht sie den Quader auf seine Grundfläche zurück; verläuft sie ausserhalb, kippt sie ihn um.

Im Quader liegt S von der unteren Kante aus $b/2$ entlang der Ebene und $h_\mathrm{S}$ senkrecht dazu. Die Verbindungslinie von der Kante zu S bildet mit der Senkrechten auf die Ebene den Winkel $\vartheta$ mit $\tan\vartheta = \frac{b/2}{h_\mathrm{S}}$. Die Senkrechte auf die Ebene ist gegenüber der Lotrechten um $\alpha$ geneigt, und zwar zur Kante hin. Ist $\alpha > \vartheta$, so liegt S jenseits der Lotrechten durch die Kante, und der Quader kippt:
\[
 \tan\alpha > \frac{b/2}{h_\mathrm{S}}
\]
Bei homogenen Quadern ist $h_\mathrm{S} = h/2$, also $\frac{b/2}{h_\mathrm{S}} = \frac{b}{h}$.

\abc
\begin{align}
 \tan\alpha &= \tan\al = \tal \approx \talP{0}{3}\\
 \text{A:}\quad \frac{b_\mathrm{A}}{h_\mathrm{A}} &= \frac{\wA}{\hA} = \qA \approx \qAP{0}{3}\\
 \text{B:}\quad \frac{b_\mathrm{B}}{h_\mathrm{B}} &= \frac{\wB}{\hB} = \qB \approx \qBP{0}{3}\\
 \text{C:}\quad \frac{b_\mathrm{C}}{h_\mathrm{C}} &= \frac{\wC}{\hC} = \qC \approx \qCP{0}{3}
\end{align}
Nur bei A ist das Verhältnis kleiner als $\tan\alpha$: \result{\text{Nur A kippt}}; B und C bleiben stehen, C nur knapp. Auf die Grösse allein kommt es nicht an: Der höchste Quader A ist auch der schmalste.

\abc
Ein Quader beginnt zu kippen bei $\tan\alpha = b/h$:
\begin{align}
 \alpha_\mathrm{A} &= \thAF = \thA \approx \result{\thAP{0}{3}}\\
 \alpha_\mathrm{B} &= \thBF = \thB \approx \result{\thBP{0}{3}}\\
 \alpha_\mathrm{C} &= \thCF = \thC \approx \result{\thCP{0}{3}}
\end{align}
Bei \alO\ ist nur der Grenzwinkel von A überschritten.
\end{abcliste}
